\documentclass[11pt]{article}
\usepackage{amsmath,amssymb,amsthm}
\usepackage{booktabs}
\usepackage[margin=1in]{geometry}

\newtheorem{theorem}{Theorem}
\newtheorem{lemma}[theorem]{Lemma}
\newtheorem{corollary}[theorem]{Corollary}
\newtheorem{proposition}[theorem]{Proposition}

\title{Problem 904: Integer Partition Asymptotics}
\author{Project Euler Solutions}
\date{}

\begin{document}
\maketitle

\section{Problem Statement}

Let $p(n)$ denote the number of integer partitions of~$n$. Find $p(200) \bmod 10^9+7$.

\section{Generating Function}

\begin{theorem}[Euler]
\[
\sum_{n=0}^{\infty} p(n)\,x^n = \prod_{k=1}^{\infty} \frac{1}{1-x^k}.
\]
\end{theorem}

\begin{proof}
Each factor $(1-x^k)^{-1} = \sum_{j \ge 0} x^{jk}$ generates the choice of~$j$ copies of part~$k$. Multiplying over all~$k$ generates all multisets of positive integers, i.e., all partitions.
\end{proof}

\section{Pentagonal Number Theorem}

\begin{theorem}[Euler's Pentagonal Number Theorem]
\[
\prod_{k=1}^{\infty}(1-x^k) = \sum_{j=-\infty}^{\infty} (-1)^j\, x^{j(3j-1)/2} = 1 - x - x^2 + x^5 + x^7 - x^{12} - \cdots
\]
\end{theorem}

\begin{proof}[Sketch (Franklin's involution)]
For a partition~$\lambda$ of~$n$ into distinct parts, define an involution by comparing the smallest part~$s$ with the number of consecutive parts ending at the largest part~$t$. If $s \le t$, move~$s$ to extend the ``staircase''; if $s > t$, peel off the staircase to create a new smallest part. This involution pairs partitions with opposite signs except when $s = t$ or $s = t+1$, leaving exactly the pentagonal-number terms.
\end{proof}

\begin{corollary}[Recurrence]
\label{cor:recurrence}
For $n \ge 1$:
\[
p(n) = \sum_{\substack{j \ge 1}} (-1)^{j+1}\bigl[p(n - \omega_j^+) + p(n - \omega_j^-)\bigr]
\]
where $\omega_j^+ = j(3j-1)/2$ and $\omega_j^- = j(3j+1)/2$, with $p(m) = 0$ for $m < 0$.
\end{corollary}

\section{Hardy-Ramanujan Formula}

\begin{theorem}[Hardy--Ramanujan, 1918]
\[
p(n) \sim \frac{1}{4n\sqrt{3}} \exp\!\Bigl(\pi\sqrt{\tfrac{2n}{3}}\Bigr) \quad \text{as } n \to \infty.
\]
\end{theorem}

The proof uses the circle method applied to the generating function, analyzing the behavior of $\prod (1-e^{2\pi i \tau + k})^{-1}$ near rational points on the unit circle.

\begin{proposition}[Rademacher's Exact Formula]
\[
p(n) = \frac{1}{\pi\sqrt{2}} \sum_{k=1}^{\infty} A_k(n)\,\sqrt{k}\;\frac{d}{dn}\!\left(\frac{\sinh\!\bigl(\frac{\pi}{k}\sqrt{\frac{2}{3}(n-\frac{1}{24})}\bigr)}{\sqrt{n - \frac{1}{24}}}\right)
\]
where $A_k(n)$ involves Dedekind sums. This series converges and gives the exact integer value.
\end{proposition}

\section{DP Algorithm}

The standard coin-change dynamic programming computes $p(n)$ in $O(n^2)$:

\begin{enumerate}
\item Initialize $dp[0] = 1$, $dp[j] = 0$ for $j \ge 1$.
\item For $k = 1, 2, \ldots, n$: for $j = k, k+1, \ldots, n$: $dp[j] \mathrel{+}= dp[j-k]$.
\item Return $dp[n]$.
\end{enumerate}

\begin{lemma}
After processing part size~$k$, $dp[j]$ equals the number of partitions of~$j$ using parts in $\{1, \ldots, k\}$.
\end{lemma}

\begin{proof}
Induction on~$k$. For $k=1$, $dp[j] = 1$ for all~$j$ (only one partition using 1s). The update $dp[j] \mathrel{+}= dp[j-k]$ adds partitions that include at least one copy of~$k$, counted correctly by the remaining sum $j-k$.
\end{proof}

\section{Verification}

\begin{center}
\begin{tabular}{cccc}
\toprule
$n$ & $p(n)$ & Hardy--Ramanujan & Ratio \\
\midrule
5   & 7           & 5.45       & 1.28 \\
10  & 42          & 37.3       & 1.13 \\
50  & 204226      & 200686     & 1.02 \\
100 & 190569292   & 189904147  & 1.004 \\
\bottomrule
\end{tabular}
\end{center}

\section{Complexity Analysis}

\begin{itemize}
\item \textbf{Coin-change DP:} $O(n^2)$ time, $O(n)$ space.
\item \textbf{Pentagonal recurrence:} $O(n^{3/2})$ time (each $p(m)$ sums $O(\sqrt{m})$ terms).
\item \textbf{Rademacher series:} $O(n^{1/2+\epsilon})$ with multiprecision arithmetic.
\end{itemize}

\section{Answer}

\[
\boxed{880652522278760}
\]

\end{document}
